\section{Алгоритмізація та програмна реалізація імітаційної моделі системи}

\subsection{Вибір засобів програмної реалізації}

Модель реалізовано мовою Java~\cite{java-jls} (версія мови~23) із системою збирання
Maven~\cite{maven}. По-перше, мовою Java написане ядро PetriObjModelPaint~\cite{petriobjmodelpaint}
(модуль \texttt{petri-math}), у якому реалізовано алгоритм імітації стохастичної мережі Петрі за
посібником~\cite{stetsenko-kursova,stetsenko-modeling}: програма підключає його як бібліотеку, тож
алгоритм імітації й семантика елементів мережі залишаються такими самими, як у ядрі. По-друге,
наслідування з перевизначенням методів дає змогу внести потрібні задачі зміни в класах-нащадках,
не змінюючи коду ядра. По-третє,
генератори з заданим seed (\texttt{SplittableRandom}~\cite{java-api}) забезпечують відтворюваність
прогонів, а скомпільована програма працює швидко: серія прогонів верифікації (понад
$4{,}8\cdot10^8$ подій) триває близько 11~с. Недолік Java -- бідніші, ніж у Python, засоби
статистичної обробки й побудови графіків, тому результати експериментів оброблено мовою Python
(розділ~4).

З ядра використано класи \texttt{PetriP}, \texttt{PetriT}, \texttt{ArcIn}, \texttt{ArcOut},
\texttt{PetriNet}, \texttt{PetriSim} з процедурами входу (\texttt{input}) і виходу маркерів
(\texttt{output}) та пошуку найближчої події (\texttt{eventMin}); вихідний код ядра не змінювався,
модифікації внесено в класах-нащадках (таблиця~\ref{tab:kernel-mods}). Ядро універсальне, тому
його використання пристосовано до задачі. Модель -- одна мережа Петрі,
тож замість циклу \texttt{PetriObjModel.go}, що на кожній події перемішує й сортує Петрі-обʼєкти
(це потрібно лише для спільних позицій кількох обʼєктів), і методу \texttt{PetriSim.go}, що друкує
маркірування на кожному кроці, працює власний цикл однієї мережі (алгоритм~\ref{alg:main}).
Решту засобів ядра задача потребує: пріоритети (стратегії 1--4), імовірності запуску (стратегія~5),
багатоканальні переходи (пари AA, AB, BB), кратні та інформаційну дуги (резервування памʼяті), тому
їх не вилучено. Натомість пошук активних переходів
\texttt{findActiveT}, на який припадає найбільше часу, перевизначено: ядро після кожного запуску
перевіряє всі переходи за списками \texttt{Integer} і сортує активні, а в \texttt{PetriSimulator}
переходи впорядковано за пріоритетом один раз до прогону й вхідні дуги зведено в масиви. Результат
той самий, а прогін у 1,25--1,28~раза швидший, ніж циклом
\texttt{PetriObjModel.go}.

\begin{kpitable}{|L{3.4cm}|L{4.8cm}|L{7.4cm}|}{Модифікації ядра PetriObjModelPaint\label{tab:kernel-mods}}
  {\textbf{Що змінено} & \textbf{В ядрі PetriObjModelPaint} & \textbf{Розширення в програмі моделі та обґрунтування}}
  Затримка переходу & метод \texttt{getTimeServ} класу \texttt{PetriT}: закон -- рядком, значення -- \texttt{FunRand}
    на основі \texttt{Math.random()} & \texttt{StochasticTransition}: затримка -- обʼєкт \texttt{Delay},
    значення -- з генератора прогону із заданим seed (прогін відтворюється точно) \\ \hline
  Вирішення конфлікту & метод \texttt{doConflikt} класу \texttt{PetriSim} на основі \texttt{Math.random()} &
    \texttt{PetriSimulator}: те саме правило (\texttt{ConflictResolver}) з генератором прогону \\ \hline
  Середні значення & \texttt{getMean} з вагою $\Delta t/t'$, тобто від моменту~0 &
    \texttt{PetriSimulator.run} викликає \texttt{doStatistics} ядра з вагою
    $\Delta t/(t'-T_{\text{п}})$: середні~\eqref{eq:time-means} лише на $[T_{\text{п}};\,T_{\text{мод}}]$ \\ \hline
  Лічильники спрацьовувань & відсутні & \texttt{StochasticTransition}: лічильники запусків і
    завершень, обнулення в момент~$T_{\text{п}}$; потрібні для $N_{\text{Arr}i}$ у формулі~\eqref{eq:L-lambda} \\ \hline
  Нумерація елементів & статичний лічильник класу, має збігатися з індексом у мережі &
    \texttt{PetriNetBuilder}: явні номери й перевірка структури мережі \\ \hline
  Спостереження за станом & один спостерігач фаз переходу & \texttt{PetriSimulator}: стан після
    кожної події -- спостерігачам \texttt{SimulationListener} (інваріанти, протокол подій) \\ \hline
  Цикл імітації і пошук активних переходів & \texttt{PetriObjModel.go} перемішує й сортує Петрі-обʼєкти
    на кожній події; \texttt{findActiveT} перевіряє всі переходи й сортує активні &
    \texttt{PetriSimulator}: цикл однієї мережі; переходи впорядковано за~$r$ один раз до прогону,
    вхідні дуги -- у масивах (той самий результат, прогін у 1,25--1,28~раза швидший) \\ \hline
\end{kpitable}

\subsection{Алгоритм імітації стохастичної мережі Петрі}

Алгоритм імітації подієвий: модельний час просувається до найближчого моменту виходу маркерів
з переходу. Стан мережі -- маркірування позицій і множина активних каналів переходів із
моментами їх завершення; кожен запуск відкриває новий канал, вилучаючи маркери з вхідних позицій
(крім інформаційних дуг), а у вихідні позиції маркери надходять через затримку~$d$
(алгоритм~\ref{alg:main}).

\begin{algorithm}[tp]
  \caption{Імітація стохастичної мережі Петрі}\label{alg:main}
  \begin{algorithmic}[1]
    \Require мережа в початковому маркіруванні, $T_{\text{мод}}$, $T_{\text{п}}$, seed генератора
    \State $t \gets 0$; якщо $T_{\text{п}} = 0$, розпочати збір статистики
    \State \Call{ВхідМаркерів}{$t$}
    \Loop
      \State $t' \gets \min_T \min_j t^{\text{вих}}_{T,j}$
      \If{$t' > T_{\text{мод}}$}
        \State \Call{ПросуванняЧасу}{$T_{\text{мод}}$}; \textbf{завершити}
      \EndIf
      \State \Call{ПросуванняЧасу}{$t'$}
      \ForAll{канали $j$ переходів $T$, де $t^{\text{вих}}_{T,j} = t$}
        \State вихід маркерів з каналу $j$: $m_P \gets m_P + k$ для дуг $(T, P, k)$
      \EndFor
      \State \Call{ВхідМаркерів}{$t$}; повідомити спостерігачів про новий стан
    \EndLoop
    \Statex
    \Procedure{ПросуванняЧасу}{$t'$}
      \If{статистика не збирається \textbf{і} $t' \ge T_{\text{п}}$}
        \State $t \gets T_{\text{п}}$; розпочати збір статистики; обнулити лічильники
      \EndIf
      \If{статистика збирається}
        \State $a \gets (t'-t)/(t'-T_{\text{п}})$
        \State для кожної позиції $P$: $\bar m_P \gets \bar m_P + (m_P-\bar m_P)\,a$
        \State для кожного переходу $T$: $\bar n_T \gets \bar n_T + (n_T-\bar n_T)\,a$
      \EndIf
      \State $t \gets t'$
    \EndProcedure
  \end{algorithmic}
\end{algorithm}

Просування стану складається з двох фаз: вихід маркерів з каналів, що завершуються в момент~$t$,
і запуск переходів, умова запуску яких виконана (алгоритм~\ref{alg:input}). Тут реалізовано
пріоритети й імовірності запуску з таблиці~\ref{tab:strategy-params}: запускається один перехід
з найвищим~$r$ (за кількох -- обраний випадково з імовірностями~$b$), після чого умови
перераховуються. Канал ResC ($d=0$) завершується в той самий момент, тож памʼять потрапляє в RC на
наступній ітерації без просування часу.

\begin{algorithm}[tp]
  \caption{Вхід маркерів у переходи з вирішенням конфліктів}\label{alg:input}
  \begin{algorithmic}[1]
    \Procedure{ВхідМаркерів}{$t$}
      \State $\mathcal{A} \gets \{T : b_T \neq 0 \text{ і } m_P \ge k \text{ для кожної вхідної дуги } (P, T, k)\}$
      \While{$\mathcal{A} \neq \emptyset$}
        \State упорядкувати $\mathcal{A}$ за спаданням $r$; $\mathcal{G}$ -- переходи з найвищим $r$
        \If{$|\mathcal{G}| = 1$}
          \State $T^* \gets$ єдиний перехід групи $\mathcal{G}$
        \Else
          \State $u \gets U[0;1)$; $T^* \gets$ перший $T \in \mathcal{G}$, для якого
            $\sum p_T > u$, де $p_T = b_T$, якщо $b_T < 1$, інакше $p_T = 1/|\mathcal{G}|$
        \EndIf
        \State $m_P \gets m_P - k$ для кожної вхідної дуги $(P, T^*, k)$, крім інформаційних
        \State відкрити канал $T^*$ з моментом виходу $t + d_{T^*}$, $d_{T^*}$ -- реалізація затримки
        \State перерахувати $\mathcal{A}$
      \EndWhile
    \EndProcedure
  \end{algorithmic}
\end{algorithm}

Між двома подіями стан не змінюється, тому рекурентна формула з вагою $a=(t'-t)/(t'-T_{\text{п}})$
дає точне середнє за часом на $[T_{\text{п}};\,t']$, а наприкінці прогону -- середні~\eqref{eq:time-means};
перший крок після $T_{\text{п}}$ має вагу $a=1$, тож середнє починається зі стану в момент~$T_{\text{п}}$. Потім
відгуки обчислюються за формулами~\eqref{eq:L-lambda}--\eqref{eq:U}.

\subsection{Структура програми}\label{sec:structure}

Програму поділено на пакети: \texttt{petri} -- розширення ядра (не залежить від задачі B14),
\texttt{model} -- модель системи B14, \texttt{verification} -- засоби верифікації,
\texttt{experiment} -- серії прогонів розділу~4, \texttt{io} і \texttt{app} -- файли CSV і запуск з
командного рядка. Залежності між пакетами спрямовані до пакета \texttt{petri}, а від нього -- до
бібліотеки ядра, тому розширення ядра перевіряється окремо від моделі. Призначення основних класів
наведено в таблиці~\ref{tab:classes}.

\begin{kpitable}{|L{5.9cm}|L{9.7cm}|}{Призначення основних класів програми\label{tab:classes}}
  {\textbf{Клас} & \textbf{Призначення}}
  \texttt{StochasticTransition} & нащадок \texttt{PetriT}: затримка з генератора прогону, лічильники \\ \hline
  \texttt{PetriSimulator} & нащадок \texttt{PetriSim}: алгоритми~\ref{alg:main} і~\ref{alg:input}, статистика після $T_{\text{п}}$ \\ \hline
  \texttt{PetriNetBuilder} & побудова й перевірка мережі \texttt{PetriNet} з елементів ядра \\ \hline
  \texttt{Strategy} & $r_A$, $r_B$, $r_C$, $k_C$, $b_A$, $b_B$ стратегій (таблиця~\ref{tab:strategy-params}) \\ \hline
  \texttt{ModelParameters} & вхідні змінні $s$, $\lambda_i$ ($\gamma\lambda_i^0$), а також $1/\mu_i$ і~$M$ \\ \hline
  \texttt{B14Net} & мережа рисунка~\ref{fig:petri}, ліва частина інваріанту~\eqref{eq:mem-invariant} \\ \hline
  \texttt{OutputCalculator} & відгуки за формулами~\eqref{eq:L-lambda}--\eqref{eq:U} \\ \hline
  \texttt{B14Simulation} & один прогін: мережа, генератор, імітація, відгуки \\ \hline
  \texttt{InvariantMonitor} & перевірка інваріанту й правил упорядкування після кожної події \\ \hline
  \texttt{VerificationRunner} & серії верифікації, зокрема порівняння з M/M/$c$ (\texttt{MmcQueue}) \\ \hline
  \texttt{ResponseTrajectory} & середні відгуки на $[0;\,t]$ і поточна черга (вибір $T_{\text{п}}$, $T_{\text{мод}}$) \\ \hline
  \texttt{ExperimentRunner} & серії прогонів експериментального дослідження \\ \hline
\end{kpitable}

Мережу будує клас \texttt{B14Net}; назви позицій і переходів збігаються з рисунком~\ref{fig:petri},
стратегію задають лише параметри переходів ResC і Serv$i$, тож структура мережі однакова для всіх
стратегій. Фрагмент побудови цих переходів
(початкова версія -- модель за умовою задачі; зміни для пропозиції -- у
підрозділі~\ref{sec:improvement}):
\begin{Verbatim}[gobble=8]
        int kC = strategy.reservationMultiplicity();
        reservationMultiplicity = kC;
        reservation = builder.transition(RESERVE_C).delay(Delay.constant(0))
                .priority(strategy.priority(TaskType.C))
                .informationalInput(queues.get(TaskType.C), 1)
                .input(memory, kC).output(reserved, kC).add();

        for (TaskType type : TaskType.values()) {
            PetriP done = builder.place("Out" + type, 0);
            PetriP resource = type == TaskType.C ? reserved : memory;
            int v = type.memoryDemand();
            int priority = type == TaskType.C ? 0 : strategy.priority(type);
            services.put(type, builder.transition(serviceName(type))
                    .delay(Delay.exponential(parameters.meanServiceTime(type)))
                    .priority(priority).probability(strategy.launchProbability(type))
                    .input(queues.get(type)).input(resource, v)
                    .output(memory, v).output(done).add());
        }
\end{Verbatim}

Звʼязки між класами, їх поля й методи показано на UML-діаграмі (рисунок~\ref{fig:uml}).

\begin{figure}
  \centering
  % Рисунок 3.1 — UML-діаграма класів програми (sim/, пакет ua.kpi.km.b14).
% Члени класів скорочено до тих, що потрібні для розуміння алгоритму:
% «+» — public, «-» — private, «~» — доступ у межах пакета.
% Пакет petri-math — бібліотека ядра (без змін); методи, показані в нащадках
% (StochasticTransition, PetriSimulator), перевизначають однойменні методи ядра.
\begin{tikzpicture}[
    font=\setstretch{1}\fontsize{8}{9.4}\selectfont,
    uml/.style={draw, fill=white, rectangle split, rectangle split parts=3,
      rectangle split ignore empty parts, rectangle split part align={center,left,left},
      text width=4.95cm, inner sep=2pt, anchor=north},
    pkg/.style={draw, dashed, rounded corners=2pt, inner sep=5pt},
    pkgname/.style={anchor=south west, inner sep=1pt,
      font=\setstretch{1}\fontsize{9}{11}\selectfont\bfseries},
    dep/.style={-{Stealth[length=2mm]}, dashed},
    assoc/.style={-{Stealth[length=2mm]}},
    comp/.style={{Diamond[length=2.6mm, width=2mm]}-{Stealth[length=2mm]}},
    agg/.style={{Diamond[open, length=2.6mm, width=2mm]}-{Stealth[length=2mm]}},
    real/.style={-{Triangle[open, length=2.4mm, width=2.4mm]}, dashed},
    inherit/.style={-{Triangle[open, length=2.4mm, width=2.4mm]}},
    ml/.style={font=\setstretch{1}\fontsize{7.5}{9}\selectfont, inner sep=1.5pt},
  ]
  \def\xa{0} \def\xb{5.7} \def\xc{11.4}
  \newcommand\tld{\textasciitilde}

  % ---------------- пакет model ----------------
  \node[uml] (Params) at (\xa,0) {\centering «record»\\\textbf{ModelParameters}
    \nodepart{second}+ strategy: Strategy\\+ arrivalRates, meanServiceTimes:\\\quad Map‹TaskType, Double›\\+ memory: int
    \nodepart{third}+ ofIntensityFactor(s, $\gamma$)\\+ ofRates($\lambda_A$, $\lambda_B$, $\lambda_C$)\\+ withArrivalRate, withMeanServiceTime,\\\quad withMemory\\+ offeredLoad(): double};
  \node[uml] (Strategy) at (\xa,-3.95) {\centering «enum»\\\textbf{Strategy}
    \nodepart{second}S1, S2, S3, S4, S5
    \nodepart{third}+ priority(TaskType): int\\+ launchProbability(TaskType): double\\+ reservationMultiplicity(): int};
  \node[uml] (TaskType) at (\xa,-6.25) {\centering «enum»\\\textbf{TaskType}
    \nodepart{second}A, B, C
    \nodepart{third}+ baseArrivalRate(): double\\+ meanServiceTime(): double\\+ memoryDemand(): int};
  \node[uml] (Settings) at (\xa,-8.55) {\centering «record»\\\textbf{RunSettings}
    \nodepart{second}+ simulationTime, warmupTime: double\\+ seed: long};

  \node[uml] (Sim) at (\xb,0) {\centering \textbf{B14Simulation}
    \nodepart{third}+ run(ModelParameters, RunSettings,\\\quad listenerFactory): SimulationResult};
  \node[uml] (Calc) at (\xb,-1.75) {\centering \textbf{OutputCalculator}
    \nodepart{third}\tld{} calculate(B14Net, ModelParameters,\\\quad RunSettings): SimulationResult};
  \node[uml] (Net) at (\xb,-3.6) {\centering \textbf{B14Net}
    \nodepart{second}- net: PetriNet\\- queues: EnumMap‹TaskType, PetriP›\\- arrivals, services: EnumMap\\- reservationLimit: PetriP (LimC)
    \nodepart{third}+ create(ModelParameters,\\\quad RandomGenerator): B14Net\\+ queue(type): PetriP\\+ service(type): StochasticTransition\\+ memoryBalance(): int};

  \node[uml] (Result) at (\xc,0) {\centering «record»\\\textbf{SimulationResult}
    \nodepart{second}+ byType: Map‹TaskType, TypeMetrics›\\+ queueLength, waitingTime,\\\quad sojournTime, utilization,\\\quad reservedIdle: double};
  \node[uml] (Metrics) at (\xc,-2.85) {\centering «record»\\\textbf{TypeMetrics}
    \nodepart{second}+ queueLength, inSystem: double\\+ arrivalRate: double\\+ waitingTime, sojournTime: double\\+ arrivals, completed: long};

  % ---------------- пакет verification ----------------
  \node[uml] (Runner) at (\xc,-5.75) {\centering \textbf{VerificationRunner}
    \nodepart{third}+ analyticTable(): CsvTable\\+ variationTable(): CsvTable\\+ singleChangeTable(): CsvTable\\+ instabilityTable(): CsvTable};
  \node[uml] (Mmc) at (\xc,-7.95) {\centering «record»\\\textbf{MmcQueue}
    \nodepart{second}+ arrivalRate, serviceRate: double\\+ servers: int
    \nodepart{third}+ queueLength(), waitingTime(),\\\quad sojournTime(): double};
  \node[uml] (Monitor) at (\xc,-10.4) {\centering \textbf{InvariantMonitor}
    \nodepart{second}- net: B14Net\\- highPriorityC: boolean
    \nodepart{third}+ onStateChanged(t)};

  % ---------------- пакет petri ----------------
  \node[uml] (Listener) at (\xc,-13.55) {\centering «interface»\\\textbf{SimulationListener}
    \nodepart{third}+ onStateChanged(t)\\+ onTransitionStarted(t, PetriT)\\+ onTransitionCompleted(t, PetriT)};
  \node[uml] (Simulator) at (\xc,-15.95) {\centering \textbf{PetriSimulator}
    \nodepart{second}- conflictResolver: ConflictResolver\\- warmupTime: double
    \nodepart{third}+ run(simulationTime, warmupTime)\\+ findActiveT(): ArrayList‹PetriT›\\+ doConflikt(enabled): PetriT\\- advanceTimeTo(time)};
  \node[uml, text width=3.9cm] (Resolver) at (\xc-0.95,-18.75) {\centering \textbf{ConflictResolver}
    \nodepart{second}- random: RandomGenerator
    \nodepart{third}+ choose(enabled): PetriT};

  \node[uml] (STrans) at (\xb,-13.55) {\centering \textbf{StochasticTransition}
    \nodepart{second}- delay: Delay\\- random: RandomGenerator\\- starts, completions: long
    \nodepart{third}+ getTimeServ(): double\\+ actIn(places, t)\\+ actOut(places, t)\\+ completions(): long};

  \node[uml] (Builder) at (\xa,-13.55) {\centering \textbf{PetriNetBuilder}
    \nodepart{third}+ place(name, m): PetriP\\+ transition(name): TransitionBuilder\\+ build(): PetriNet};
  \node[uml] (Delay) at (\xa,-15.75) {\centering «sealed interface»\\\textbf{Delay}
    \nodepart{second}Constant(value), Exponential(mean)
    \nodepart{third}+ sample(RandomGenerator): double};

  % ---------------- пакет experiment (розділ 4) ----------------
  \node[uml, text width=3.3cm] (Traj) at (-0.75,-10.9) {\centering \textbf{ResponseTrajectory}
    \nodepart{second}- net: B14Net\\- step: double
    \nodepart{third}+ onStateChanged(t)\\+ points(): List‹Point›};
  \node[uml, text width=3.3cm] (Exp) at (2.85,-10.9) {\centering \textbf{ExperimentRunner}
    \nodepart{third}+ run(series, out, \ldots)\\+ factorialTable(\ldots)\\\tld{} seed(series, cell, k)};

  % ---------------- бібліотека petri-math (ядро) ----------------
  \node[uml] (KNet) at (\xa,-20.75) {\centering \textbf{PetriNet}
    \nodepart{second}- ListP: PetriP[]\\- ListT: PetriT[]\\- ListIn: ArcIn[]; ListOut: ArcOut[]};
  \node[uml] (KT) at (\xb,-20.75) {\centering \textbf{PetriT}
    \nodepart{second}- priority: int; probability: double\\- buffer: int; timeOut: List‹Double›
    \nodepart{third}+ actIn(places, t); actOut(places, t)\\+ getTimeServ(); getMean(): double};
  \node[uml] (KSim) at (\xc,-20.75) {\centering \textbf{PetriSim}
    \nodepart{third}+ findActiveT(): ArrayList‹PetriT›\\+ input(); output(); eventMin()\\+ doConflikt(list): PetriT\\+ doStatistics(a)};

  \coordinate (k2) at ($(KT.south)+(0,-0.3)$); % траса асоціації PetriSim -> PetriNet

  % ---------------- рамки пакетів ----------------
  \begin{scope}[on background layer]
    \node[pkg, fit=(Params)(Settings)(Sim)(Net)(Result)(Metrics)] (Pmodel) {};
    \node[pkg, fit=(Runner)(Mmc)(Monitor)] (Pver) {};
    \node[pkg, fit=(Builder)(Delay)(STrans)(Listener)(Resolver)] (Ppetri) {};
    \node[pkg, fit=(KNet)(KT)(KSim)(k2)] (Pkernel) {};
    \node[pkg, fit=(Traj)(Exp)] (Pexp) {};
  \end{scope}
  \node[pkgname] at (Pmodel.north west) {model};
  \node[pkgname] at (Pver.north west) {verification};
  \node[pkgname] at (Ppetri.north west) {petri};
  \node[pkgname] at (Pkernel.north west) {petri-math (ядро PetriObjModelPaint)};
  \node[pkgname] at (Pexp.north west) {experiment};

  % ---------------- зв'язки ----------------
  % model
  \draw[assoc] (Params) -- (Strategy);
  \draw[assoc] (Strategy) -- (TaskType);
  \coordinate (s1) at ([yshift=2mm]Sim.west);
  \draw[dep] (s1) -- (Params.east |- s1);
  \coordinate (g1) at ($([yshift=-2mm]Sim.west)+(-0.3,0)$);
  \draw[dep] ([yshift=-2mm]Sim.west) -- (g1) |- (Net.west);
  \draw[dep] (g1) |- (Settings.east);
  \draw[dep] (Sim) -- (Calc);
  \draw[dep] (Calc) -- (Net);
  \draw[dep] (Sim.east) -- (Result.west |- Sim.east);
  \draw[dep] (Calc.east) -| ($(Result.south west)+(0.5,0)$);
  \draw[comp] (Result) -- node[ml, right]{3} (Metrics);
  % verification
  \draw[dep] (Runner) -- (Mmc);
  \draw[dep] ([yshift=-2mm]Runner.west) -- ++(-0.32,0) |- ([yshift=2mm]Monitor.west);
  \draw[assoc] ([yshift=-2mm]Monitor.west) -- ++(-0.5,0) |- (Net.east);
  \draw[real] (Monitor) -- (Listener);
  % petri
  \draw[assoc] ([xshift=0.8cm]Net.south) -- node[ml, right, pos=0.93]{*} ([xshift=0.8cm]Net.south |- STrans.north);
  \coordinate (b1) at ($([xshift=-0.8cm]Net.south)+(0,-6.27)$);
  \draw[dep] ([xshift=-0.8cm]Net.south) -- (b1) -| (Builder.north);
  \draw[dep] (Builder.east) -- (Builder.east -| STrans.west);
  \draw[assoc] (STrans.west |- Delay.east) -- (Delay.east);
  \draw[agg] (Simulator) -- node[ml, right]{*} (Listener);
  \draw[comp] (Resolver.north |- Simulator.south) -- (Resolver.north);
  % experiment: залежність від пакета model (B14Simulation, ModelParameters, RunSettings)
  \draw[dep] (Exp.west |- Traj.east) -- (Traj.east);
  \draw[dep] ([xshift=1.0cm]Exp.north) -- ([xshift=1.0cm]Exp.north |- Pmodel.south);
  \draw[real] (Traj.south) -- ++(0,-0.3) -| ([xshift=-1.6cm]Listener.north);
  % наслідування від класів ядра
  \draw[inherit] (STrans) -- (KT);
  \draw[inherit] ([xshift=2.0cm]Simulator.south) -- ([xshift=2.0cm]Simulator.south |- KSim.north);
  % petri-math
  \coordinate (k1) at ($(Builder.west)+(-0.28,0)$);
  \draw[dep] (Builder.west) -- (k1) |- (KNet.west);
  \draw[comp] (KNet.east |- KT.west) -- node[ml, above]{*} (KT.west);
  \draw[assoc] (KSim.south) -- (KSim.south |- k2) -| ([xshift=1.4cm]KNet.south);
\end{tikzpicture}

  \caption{UML-діаграма класів програми}
  \label{fig:uml}
\end{figure}

Клас \texttt{StochasticTransition} генерує затримку законом \texttt{Delay} і оновлює лічильники
після виклику методу ядра, лише якщо кількість активних каналів змінилася:
\codeinput[firstline=46, lastline=72, gobble=4]{../sim/src/main/java/ua/kpi/km/b14/petri/StochasticTransition.java}

У методі просування часу класу \texttt{PetriSimulator} в момент~$T_{\text{п}}$ обнуляються
лічильники, а вага методу ядра \texttt{doStatistics} відлічується від~$T_{\text{п}}$:
\codeinput[firstline=124, lastline=152, gobble=4]{../sim/src/main/java/ua/kpi/km/b14/petri/PetriSimulator.java}

Клас \texttt{ConflictResolver} відтворює правило ядра: перехід з $b=1$ у групі з $n$
рівнопріоритетних отримує ймовірність $1/n$, інший -- свою~$b$ (у стратегії~5 -- по~0,5):
\codeinput[firstline=36, lastline=52, gobble=4]{../sim/src/main/java/ua/kpi/km/b14/petri/ConflictResolver.java}

Клас \texttt{OutputCalculator} обчислює відгуки за формулами~\eqref{eq:L-lambda}
і~\eqref{eq:W-type}:
\codeinput[firstline=28, lastline=44, gobble=8]{../sim/src/main/java/ua/kpi/km/b14/model/OutputCalculator.java}

Програму запускають з командного рядка (команди \texttt{run}, \texttt{verify},
\texttt{experiment}). Програму перевірено 42~тестами JUnit~5~\cite{junit}, головні -- тести
еквівалентності розширення ядра: мережа з пріоритетами, кратними й інформаційною дугами та багатоканальними
переходами дає ті самі маркірування й середні в циклі \texttt{PetriObjModel.go} ядра і в
\texttt{PetriSimulator}, а \texttt{findActiveT} у мережі B14 -- той самий список, що й ядро. Інші
тести перевіряють точні лічильники й середні для мережі з детермінованими затримками, початок
статистики після~$T_{\text{п}}$, перевагу пріоритету, невилучення маркерів інформаційною дугою,
частоти вибору за ймовірностями запуску, відтворюваність прогону з тим самим seed і нумерацію
елементів незалежно від лічильників ядра. Текст
основних класів наведено в додатку~А (ядра -- у~\cite{petriobjmodelpaint}), засобів верифікації --
у додатку~В. Повний вихідний код програми, скриптів аналізу, мережу Петрі у форматі PNML і
результати прогонів (CSV) розміщено в репозиторії GitHub \repourl.

\subsection{Верифікація моделі}\label{sec:verification}

Верифікація має підтвердити, що програма відтворює систему, описану в розділах 1 і~2. Застосовано
перевірку інваріантів і правил упорядкування в кожному стані, порівняння з аналітичними
результатами для окремих випадків і аналіз зміни відгуків за зміни вхідних змінних. Seed
генератора -- \texttt{20261118} (за кількох прогонів у точці -- \texttt{20261118}…\texttt{20261127});
незмінні параметри -- за таблицею~\ref{tab:params}: $\mu_A = 0{,}25$, $\mu_B = 0{,}2$,
$\mu_C = 0{,}1$~завд./хв, $M = 2$, $v_A = v_B = 1$, $v_C = 2$ ($v_i$ -- в усіх серіях, решта -- крім
серії зі зміною кожної змінної).

\textbf{Інваріанти та правила упорядкування.} Спостерігач \texttt{InvariantMonitor} після
кожної події всіх прогонів верифікації перевіряв: інваріант~\eqref{eq:mem-invariant}; відсутність
простою памʼяті, коли в черзі є A чи~B ($m_{\text{QA}} + m_{\text{QB}} > 0 \Rightarrow
m_{\text{Mem}} = 0$); для стратегій 3--5 -- резервування всієї вільної памʼяті, поки C очікує
($m_{\text{QC}} > 0 \Rightarrow m_{\text{Mem}} = 0$); для стратегій 1--2 -- неможливість очікування~C
за двох вільних одиниць ($m_{\text{QC}} > 0 \Rightarrow m_{\text{Mem}} \le 1$). За
$4{,}8\cdot10^8$ подій порушень не виявлено: заборонених умовою станів немає, а резервування працює
так, як описано в підрозділі~\ref{sec:strategies-net}.

\textbf{Протокол подій.} Просування часу й стану перевірено за протоколом подій (додаток~Б):
переходи, з яких вийшли й у які увійшли маркери, і стан після кожної події. У протоколі стратегії~3
($\gamma = 0{,}6$) в момент 34,07~хв завершилося завдання~C за 5 завдань A і 3 завдання B у черзі:
обидві одиниці памʼяті зайняли B вищого пріоритету, а A стартували лише після спорожнення черги~B.
У момент 87,86~хв надійшло~C, коли виконувалося B: ResC зарезервував вільну одиницю
($m_{\text{RC}} = 1$), тож B, що надійшло о 90,91~хв, чекало, хоча формально памʼять не була
зайнята виконанням; після завершення B о 91,84~хв ResC
зарезервував другу одиницю і стартувало~C. У стратегії~1 C стартує лише за порожніх черг A і~B.

\textbf{Контрольні випадки з аналітичним розвʼязком.} Якщо залишити потік лише одного типу,
система стає M/M/1 для завдань~C (займають усю памʼять) і M/M/2 для A або~B, характеристики
яких дає формула Ерланга~C~\cite{kleinrock}:
\begin{equation}\label{eq:erlang}
  L_q = \frac{P_0\,a^c\,\rho}{c!\,(1-\rho)^2}, \quad
  P_0 = \left[\sum_{n=0}^{c-1}\frac{a^n}{n!} + \frac{a^c}{c!\,(1-\rho)}\right]^{-1}, \quad
  W_q = \frac{L_q}{\lambda}, \quad W = W_q + \frac{1}{\mu},
\end{equation}
де $a = \lambda/\mu$ -- навантаження, $c$ -- кількість каналів, $\rho = a/c$. Для кожного випадку
виконано 10~прогонів і побудовано довірчий інтервал рівня~0,95 ($\bar x \pm
t_{0{,}975;\,9}\,s/\sqrt{10}$): усі 12 теоретичних значень потрапили в інтервали, відхилення
не перевищує 2~\% (таблиця~\ref{tab:verif-analytic}); випадок $\rho = 0{,}8$ перевіряє режим
довгих черг.

\begin{kpitable}{|L{3.6cm}|r|r|r|r|}{Порівняння результатів імітації з аналітичними
  ($T_{\text{мод}} = 10^6$~хв, $T_{\text{п}} = 10^4$~хв, 10 прогонів)\label{tab:verif-analytic}}
  {\textbf{Відгук} & \textbf{Теорія} & \textbf{Модель} & \textbf{$\pm$ (0,95)} & \textbf{$\delta$, \%}}
  \multicolumn{5}{|l|}{лише C, $\lambda_C = 0{,}06$~завд./хв: M/M/1, $\rho = 0{,}6$} \\ \hline
  $L_{q,C}$, завд. & 0,9000 & 0,9056 & 0,0168 & 0,62 \\ \hline
  $W_{q,C}$, хв & 15,0000 & 15,1063 & 0,2878 & 0,71 \\ \hline
  $W_C$, хв & 25,0000 & 25,1026 & 0,3093 & 0,41 \\ \hline
  \multicolumn{5}{|l|}{лише A, $\lambda_A = 0{,}2$~завд./хв: M/M/2, $\rho = 0{,}4$} \\ \hline
  $L_{q,A}$, завд. & 0,1524 & 0,1516 & 0,0018 & 0,52 \\ \hline
  $W_{q,A}$, хв & 0,7619 & 0,7575 & 0,0081 & 0,58 \\ \hline
  $W_A$, хв & 4,7619 & 4,7536 & 0,0142 & 0,17 \\ \hline
  \multicolumn{5}{|l|}{лише A, $\lambda_A = 0{,}4$~завд./хв: M/M/2, $\rho = 0{,}8$} \\ \hline
  $L_{q,A}$, завд. & 2,8444 & 2,8312 & 0,0571 & 0,46 \\ \hline
  $W_{q,A}$, хв & 7,1111 & 7,0746 & 0,1400 & 0,51 \\ \hline
  $W_A$, хв & 11,1111 & 11,0723 & 0,1438 & 0,35 \\ \hline
  \multicolumn{5}{|l|}{лише B, $\lambda_B = 0{,}07$~завд./хв: M/M/2, $\rho = 0{,}175$} \\ \hline
  $L_{q,B}$, завд. & 0,0111 & 0,0109 & 0,0003 & 1,82 \\ \hline
  $W_{q,B}$, хв & 0,1580 & 0,1551 & 0,0045 & 1,84 \\ \hline
  $W_B$, хв & 5,1580 & 5,1536 & 0,0176 & 0,08 \\ \hline
\end{kpitable}

\textbf{Зміна кожної вхідної змінної та параметра окремо.} Базовий рядок відповідає $\gamma = 0{,}4$
($\lambda_A = 0{,}08$, $\lambda_B = 0{,}028$, $\lambda_C = 0{,}024$~завд./хв, $1/\mu_i = 4$; 5;
10~хв, $M = 2$, $\rho = 0{,}47$); далі подвоєно одну вхідну змінну $\lambda_i$ або
параметр $1/\mu_i$, в останньому рядку $M = 3$; скрізь $\rho \le 0{,}71$. Стратегії 1 і~3 мають
протилежне упорядкування; у точці -- 10~прогонів з
$T_{\text{мод}} = 3{,}5\cdot10^6$~хв і $T_{\text{п}} = 2\cdot10^4$~хв, як у розділі~4
(таблиця~\ref{tab:verif-single}).

\begin{kpitable}{|l|r|r|r|r|r|r|r|r|r|}{Відгуки моделі за зміни однієї вхідної змінної або параметра ($T_{\text{мод}} = 3{,}5\cdot10^6$~хв, $T_{\text{п}} = 2\cdot10^4$~хв, 10 прогонів)\label{tab:verif-single}}
  {\textbf{Змінена величина} & \textbf{$\rho$} & \textbf{$U$} & \textbf{$L_q$,} & \textbf{$W_q$,} &
   \textbf{$W$,} & \textbf{$W_{q,A}$,} & \textbf{$W_{q,B}$,} & \textbf{$W_{q,C}$,} & \textbf{$W_C$,} \\
   & & & \textbf{завд.} & \textbf{хв} & \textbf{хв} & \textbf{хв} & \textbf{хв} & \textbf{хв} & \textbf{хв}}
  \multicolumn{10}{|l|}{стратегія 1 (A~$>$~B~$>$~C)} \\ \hline
  базовий & 0,470 & 0,470 & 0,58 & 4,40 & 9,70 & 2,88 & 3,55 & 10,47 & 20,46 \\ \hline
  $\lambda_A$ $\times\,2$ & 0,630 & 0,630 & 1,31 & 6,20 & 11,01 & 3,91 & 6,16 & 21,51 & 31,52 \\ \hline
  $\lambda_B$ $\times\,2$ & 0,540 & 0,540 & 0,82 & 5,12 & 10,37 & 3,05 & 4,08 & 14,44 & 24,44 \\ \hline
  $\lambda_C$ $\times\,2$ & 0,710 & 0,710 & 2,40 & 15,36 & 21,39 & 5,54 & 6,79 & 36,75 & 46,76 \\ \hline
  $1/\mu_A$ $\times\,2$ & 0,630 & 0,629 & 1,23 & 9,32 & 17,04 & 4,24 & 6,64 & 29,45 & 39,44 \\ \hline
  $1/\mu_B$ $\times\,2$ & 0,540 & 0,540 & 0,81 & 6,11 & 12,47 & 3,23 & 4,23 & 17,90 & 27,90 \\ \hline
  $1/\mu_C$ $\times\,2$ & 0,710 & 0,710 & 2,69 & 20,40 & 27,52 & 11,10 & 13,99 & 58,87 & 78,86 \\ \hline
  $M = 3$ & 0,313 & 0,313 & 0,13 & 0,96 & 6,26 & 0,35 & 0,44 & 3,62 & 13,63 \\ \hline
  \multicolumn{10}{|l|}{стратегія 3 (C~$>$~B~$>$~A)} \\ \hline
  базовий & 0,470 & 0,471 & 0,92 & 6,94 & 12,24 & 7,84 & 5,65 & 5,44 & 15,44 \\ \hline
  $\lambda_A$ $\times\,2$ & 0,630 & 0,629 & 2,51 & 11,86 & 16,66 & 13,54 & 6,68 & 6,69 & 16,68 \\ \hline
  $\lambda_B$ $\times\,2$ & 0,540 & 0,540 & 1,42 & 8,85 & 14,10 & 11,12 & 6,76 & 6,17 & 16,17 \\ \hline
  $\lambda_C$ $\times\,2$ & 0,710 & 0,710 & 4,89 & 31,36 & 37,38 & 45,17 & 24,28 & 12,47 & 22,46 \\ \hline
  $1/\mu_A$ $\times\,2$ & 0,630 & 0,630 & 1,93 & 14,62 & 22,35 & 18,08 & 9,05 & 9,59 & 19,60 \\ \hline
  $1/\mu_B$ $\times\,2$ & 0,540 & 0,540 & 1,48 & 11,24 & 17,59 & 13,39 & 7,99 & 7,85 & 17,87 \\ \hline
  $1/\mu_C$ $\times\,2$ & 0,710 & 0,710 & 8,18 & 62,00 & 69,12 & 80,48 & 43,94 & 21,49 & 41,50 \\ \hline
  $M = 3$ & 0,313 & 0,313 & 0,13 & 1,02 & 6,32 & 0,53 & 0,39 & 3,37 & 13,36 \\ \hline
\end{kpitable}


Зміна кожної змінної та параметра дає очікуваний результат:
\begin{itemize}
  \item збільшення будь-якої $\lambda_i$ чи $1/\mu_i$ підвищує~$\rho$ і збільшує $L_q$, $W_q$, $W$, а
    збільшення памʼяті зменшує їх; $|U - \rho| \le 0{,}0012$ в усіх рядках;
  \item додаткове навантаження найбільше подовжує очікування найнижчого пріоритету: за стратегії~1
    подвоєння $\lambda_A$ збільшує $W_{q,A}$ на 36~\%, а $W_{q,C}$ -- у 2,1~раза; за стратегії~3
    подвоєння $\lambda_C$ збільшує $W_{q,C}$ у 2,3, а $W_{q,A}$ -- у 5,8~раза, бо кожне нове
    завдання~C резервує памʼять і відкладає старт завдань A і~B;
  \item подвоєння $1/\mu_C$ і~$\lambda_C$ дає те саме~$\rho$ і близьку чергу ($L_q = 2{,}69$
    і~2,40~завд. за стратегії~1), але довше очікування ($W_q = 20{,}40$ і 15,36~хв), бо
    $W_q = L_q/\lambda$, а потік не зріс -- як і в M/M/$c$, де $L_q$~\eqref{eq:erlang} залежить лише
    від~$a$;
  \item $W - W_q$ дорівнює середньому часу виконання, зваженому за інтенсивностями (7,12~хв за
    $1/\mu_C = 20$~хв);
  \item за $M = 3$ C вже не монополізує памʼять: $W_q$ менший у 4,6--6,8~раза, а стратегії майже
    рівноцінні ($W = 6{,}26$ і 6,32~хв).
\end{itemize}

\textbf{Зміна стратегії та коефіцієнта інтенсивності.} Для пʼяти стратегій і
$\gamma \in \{0{,}4;\,0{,}5;\,0{,}6;\,0{,}7\}$ виконано по 10~прогонів з тими самими параметрами
($\gamma \ge 0{,}8$ не включено, бо стратегії 1--2 втрачають стійкість, підрозділ~\ref{sec:exp-plan}).
Повні результати (20~точок) -- у таблиці~\ref{tab:verif-variation-full} додатка~В, найпоказовіші --
у таблиці~\ref{tab:verif-variation} ($\bar m_{\text{RC}}/M$ -- частка памʼяті, що простоює через
резервування). Половина довірчого інтервалу для~$W$ не перевищує 2,4~\%, а відхилення від
результатів розділу~4 (60~прогонів) не перевищує 1,5~\%.

\begin{kpitable}{|c|c|r|r|r|r|r|r|r|r|r|}{Відгуки моделі за зміни стратегії та коефіцієнта інтенсивності, вибрані точки ($T_{\text{мод}} = 3{,}5\cdot10^6$~хв, $T_{\text{п}} = 2\cdot10^4$~хв, 10 прогонів)\label{tab:verif-variation}}
  {\textbf{$s$} & \textbf{$\gamma$} & \textbf{$L_q$,} & \textbf{$W_q$,} & \textbf{$W$,} &
   \textbf{$W_{q,A}$,} & \textbf{$W_{q,B}$,} & \textbf{$W_{q,C}$,} & \textbf{$U$} & \textbf{$\rho$} &
   \textbf{$\bar m_{\text{RC}}/M$} \\
   & & \textbf{завд.} & \textbf{хв} & \textbf{хв} & \textbf{хв} & \textbf{хв} & \textbf{хв} & & & }
  1 & 0,4 & 0,58 & 4,40 & 9,70 & 2,88 & 3,55 & 10,47 & 0,470 & 0,470 & 0,000 \\ \hline
  1 & 0,7 & 5,21 & 22,55 & 27,86 & 6,10 & 9,69 & 92,34 & 0,823 & 0,823 & 0,000 \\ \hline
  2 & 0,7 & 5,05 & 21,86 & 27,16 & 8,01 & 5,13 & 87,57 & 0,822 & 0,823 & 0,000 \\ \hline
  3 & 0,7 & 13,24 & 57,27 & 62,58 & 83,29 & 21,69 & 12,03 & 0,824 & 0,823 & 0,035 \\ \hline
  4 & 0,7 & 11,68 & 50,54 & 55,84 & 33,99 & 130,93 & 12,01 & 0,822 & 0,823 & 0,035 \\ \hline
  5 & 0,7 & 13,28 & 57,45 & 62,75 & 79,90 & 32,21 & 12,07 & 0,823 & 0,823 & 0,035 \\ \hline
\end{kpitable}


Результати узгоджуються з логікою функціонування системи:
\begin{itemize}
  \item зі зростанням $\gamma$ монотонно й дедалі швидше зростають $L_q$, $W_q$, $W$ і очікування
    кожного типу; $|U - \rho| \le 0{,}0012$, тобто модель виконує весь потік;
  \item порядок часів очікування типів відповідає упорядкуванню: $W_{q,A} < W_{q,B} < W_{q,C}$ для
    стратегії~1, $W_{q,B} < W_{q,A} < W_{q,C}$ для~2, $W_{q,C} < W_{q,B} < W_{q,A}$ для~3,
    $W_{q,C} < W_{q,A} < W_{q,B}$ для~4, найменший $W_{q,C}$ для~5;
  \item резервування ($\bar m_{\text{RC}} > 0$) виникає лише в стратегіях 3--5;
  \item у стратегії~5 A чекають довше за B, хоча пріоритети рівні: за рівних імовірностей вибору
    черга інтенсивнішого потоку A обслуговується з тією самою частотою, що й черга~B, тож кожне
    завдання~A чекає довше;
  \item за $\gamma = 0{,}7$ серед стратегій 3--5 найкраща~4, а стратегії 1--2 дають у 2,0--2,3~раза
    менший~$W$, як і в розділі~4.
\end{itemize}

\textbf{Нестійкий режим за умови задачі.} За $\gamma = 1$ ($\rho = 1{,}175$) середня черга,
обчислена від моменту~0, зростає майже пропорційно часу моделювання (таблиця~\ref{tab:verif-unstable},
один прогін для кожного $T_{\text{мод}}$): за стратегії~1 -- черга завдань~C найнижчого пріоритету,
за стратегії~3 -- черга завдань~A ($W$ не наведено: закон Літтла тут не діє).

\begin{kpitable}{|c|r|r|r|r|r|}{Відгуки моделі за $\gamma = 1$ залежно від часу моделювання
  ($T_{\text{п}} = 0$, 1~прогін)\label{tab:verif-unstable}}
  {\textbf{$s$} & \textbf{$T_{\text{мод}}$, хв} & \textbf{$L_{q,A}$, завд.} & \textbf{$L_{q,B}$, завд.} &
   \textbf{$L_{q,C}$, завд.} & \textbf{$L_q$, завд.}}
  1 & 10\,000 & 1,45 & 1,04 & 90,03 & 92,5 \\ \hline
  1 & 20\,000 & 1,54 & 1,11 & 207,59 & 210,2 \\ \hline
  1 & 50\,000 & 1,51 & 1,17 & 569,75 & 572,4 \\ \hline
  1 & 100\,000 & 1,50 & 1,19 & 1186,94 & 1189,6 \\ \hline
  3 & 10\,000 & 566,32 & 4,52 & 1,13 & 572,0 \\ \hline
  3 & 20\,000 & 1089,87 & 6,69 & 1,19 & 1097,8 \\ \hline
  3 & 50\,000 & 2637,02 & 5,28 & 1,09 & 2643,4 \\ \hline
  3 & 100\,000 & 5049,73 & 5,33 & 1,14 & 5056,2 \\ \hline
\end{kpitable}

Отже, модель коректно відтворює обмеження на одночасне виконання, монополізацію системи
завданням~C і всі стратегії, а для випадків з відомим розвʼязком статистично не відрізняється від
аналітичних результатів, тобто адекватна концептуальній і формалізованій моделям.

\subsection{Висновки до розділу}

Розроблено алгоритм імітації стохастичної мережі Петрі з подієвим просуванням часу, вирішенням
конфліктів за пріоритетом і ймовірністю запуску та збором статистики після перехідного періоду.
Модель реалізовано мовою Java на основі ядра PetriObjModelPaint: сім модифікацій, зокрема
відтворюваність прогонів і усунення впливу перехідного періоду, внесено в класах-нащадках. Верифікація показала відсутність порушень інваріанту
й правил упорядкування за $4{,}8\cdot10^8$ подій, збіг з аналітичними результатами M/M/1 і M/M/2
(відхилення до 2~\%) і правильний напрям зміни відгуків за зміни кожної вхідної змінної та
параметра; за $\gamma = 1$ система нестійка.
